\section{Test}
\label{sec:test}

\subsection{Come sono fatti i test}

I test sono scritti con \code{unittest} (PyUnit), si trovano nella cartella \code{tests} e si eseguono con

\par\noindent\begin{minipage}{\linewidth}
\begin{lstlisting}
python -m unittest discover -s tests -v
\end{lstlisting}
\end{minipage}\par

Sono \textbf{65 test} e passano tutti.

Ogni test parte da dati nuovi, creati nel metodo \code{setUp}, quindi i test non dipendono l'uno dall'altro. Stati, strategie di prezzo e modelli sono provati da soli (test di unità). Il controller è provato insieme al DAO tenuto in memoria, senza file (test di integrazione); al posto del servizio di pagamento vero si usa un oggetto di prova, \code{FixedPaymentService}, che risponde sempre allo stesso modo (\emph{stub}). I casi sono scelti a partire dai requisiti, provando sia i dati giusti sia quelli sbagliati (posti liberi, bloccati, venduti, inesistenti; carte valide e scadute) e i valori al limite (timeout appena sotto e appena sopra i 15 minuti, sconto dello 0\% e del 100\%).

{\small
\begin{tabularx}{\textwidth}{@{}L{6.0cm}Y@{}}
\toprule
\textbf{File} & \textbf{Cosa verifica} \\
\midrule
\code{test_ticket_state.py} & Passaggi di stato del biglietto e operazioni non permesse. \\
\code{test_pricing_strategy.py} & Prezzo normale, sconto di gruppo, abbonamento, arrotondamento al centesimo. \\
\code{test_ticket_seller_controller.py} & Acquisto approvato e rifiutato, posti duplicati, timeout, gestione degli eventi. \\
\code{test_requirements.py} & Un test per ogni requisito e i casi corretti durante lo sviluppo. \\
\code{test_subscription.py} & Ricerca dell'abbonamento per codice e acquisto scontato. \\
\code{test_ticket_db.py} & Salvataggio e lettura dal file, mappa dei posti, vincoli sugli eventi. \\
\code{test_payment_service.py}, \code{test_model_validation.py}, \code{test_utils.py} & Controllo dei dati della carta, dei modelli e delle date. \\
\bottomrule
\end{tabularx}}

\subsection{Esempi}

Pagamento approvato: il biglietto risulta venduto.

\par\noindent\begin{minipage}{\linewidth}
\begin{lstlisting}
def test_successful_purchase_sells_ticket(self) -> None:
    controller = TicketSeller(self.db, self.db, FixedPaymentService(True), StandardPricing())
    result = controller.purchase_access_title(EventId(1), [SeatId("S001")], self.customer, self.payment)
    ticket = self.db.find_tickets_by_seats(EventId(1), [SeatId("S001")])[0]
    self.assertTrue(result.approved)
    self.assertEqual("sold", ticket.state.name())
\end{lstlisting}
\end{minipage}\par

Carta non valida: viene segnalato l'errore e il posto torna disponibile.

\par\noindent\begin{minipage}{\linewidth}
\begin{lstlisting}
def test_rf8_invalid_card_releases_reserved_seats(self) -> None:
    controller = TicketSeller(self.db, self.db, CreditCardCharges(success_rate=1.0), StandardPricing())
    bad_card = PaymentData("Alice Buyer", "not-a-card", 12, 2099, "123")
    with self.assertRaises(ValueError):
        controller.purchase_access_title(EventId(1), [SeatId("S001")], self.customer, bad_card)
    self.assertEqual("available", self._ticket("S001").state.name())
\end{lstlisting}
\end{minipage}\par

\subsection{Requisiti e test}

\small
\begin{center}
\begin{tabular}{@{}L{1.8cm}L{12.3cm}@{}}
\toprule
\textbf{Requisito} & \textbf{Test principale} \\
\midrule
RF1 & \code{test_rf1_search_future_events_excludes_past_events} \\
RF2 & \code{test_rf2_seating_map_lists_every_seat_with_state} \\
RF3 & \code{test_rf3_check_availability_reports_locked_seats}, \code{test_rf4_reserve_locks_seats_for_customer} \\
RF4 & \code{test_invalid_card_number_is_rejected} \\
RF5 & \code{test_rf5_runtime_strategy_change_applies_group_discount} \\
RF6 & \code{test_rf6_payment_service_receives_computed_total} \\
RF7 & \code{test_successful_purchase_sells_ticket} \\
RF8 & \code{test_declined_purchase_releases_ticket}, \code{test_rf8_invalid_card_releases_reserved_seats} \\
RF9, RNF7 & \code{test_rf9_rnf7_timeout_boundary} \\
RF10 & \code{test_rf10_clerk_creates_event_with_seating_map} \\
RF11 & \code{test_rf11_clerk_updates_event} \\
RF12 & \code{test_rf12_clerk_deletes_event}, \code{test_delete_event_with_reserved_ticket_is_rejected} \\
RF13 & \code{test_rf13_subscription_purchase_charges_discounted_amount} \\
RNF3 & \code{test_rnf3_concurrent_reservations_allow_one_winner} \\
RNF4 & \code{test_money_is_rounded_half_up} \\
RNF5 & \code{test_rnf5_json_roundtrip_preserves_states_owner_and_lock_time} \\
\bottomrule
\end{tabular}
\end{center}
\normalsize

\section*{Conclusioni}
\addcontentsline{toc}{section}{Conclusioni}

Il progetto segue tutte le fasi viste nel corso, dai requisiti ai test. L'architettura MVC tiene separate le finestre, la logica e i dati; il DAO nasconde il modo in cui i dati sono salvati; i pattern State e Strategy rendono chiare le regole sui biglietti e sui prezzi.

Il prototipo ha alcuni limiti: il pagamento è simulato, il cassiere non deve autenticarsi e si assume che sull'archivio lavori una postazione alla volta. Per usarlo con più postazioni basterebbe sostituire il file con un database, cambiando solo la classe \code{TicketDB}.
